% Generated draft: requires mathematical and source review.
\documentclass[11pt]{article}
\usepackage[margin=25mm]{geometry}
\usepackage{fontspec}
\setmainfont{Latin Modern Roman}
\usepackage{amsmath,amssymb,mathtools,unicode-math}
\setmathfont{Latin Modern Math}
\usepackage{enumitem,needspace,xurl,hyperref}
\setlength{\parindent}{0pt}
\setlength{\parskip}{5pt}
\setlength{\emergencystretch}{3em}
\newcommand{\N}{\mathbb{N}}
\newcommand{\Z}{\mathbb{Z}}
\newcommand{\Q}{\mathbb{Q}}
\newcommand{\R}{\mathbb{R}}
\newcommand{\C}{\mathbb{C}}
\newcommand{\F}{\mathbb{F}}
\newcommand{\A}{\mathbb{A}}
\newcommand{\PP}{\mathbb{P}}
\newcommand{\E}{\mathbb{E}}
\newcommand{\eps}{\varepsilon}
\newcommand{\dd}{\,\mathrm{d}}
\begin{document}
\section*{UnsolvedMath ID: 20001995}
\textbf{Source problem code:} catalogue entry
\subsection{Definitions and mathematical statement}
Yong-Geun Oh's item A.13 asks for a higher-genus analogue of the Fukaya--Oh degeneration of holomorphic polygons in a cotangent bundle to Morse flow graphs. The source reports that, for discs with boundary on a suitable tuple of Lagrangian graphs, the holomorphic-polygon moduli space is diffeomorphic to the Morse graph-flow moduli space in the adiabatic, or large-complex-structure, limit. Projections to the base have amoeba-like shapes that collapse to flow graphs. Precise genericity, compactification and analytic hypotheses of the disc theorem are not reproduced.

The paired aid gives necessary exactness and topology conditions. In \(T^*Q\) use the tautological form \(\lambda\), \(\omega=d\lambda\), compatible almost complex structures, and fixed graphs
\[
 L_i^\epsilon=\operatorname{graph}(\epsilon\,df_i),\qquad
 \lambda|_{L_i^\epsilon}=\epsilon\,df_i,\qquad\epsilon>0.
\]
Consider unperturbed finite-energy pseudoholomorphic maps from punctured bordered surfaces, with boundary arcs labeled by these graphs and the asymptotic/corner convergence needed for Stokes' theorem. For an oriented boundary arc \(a\), let \(q_a^-\) and \(q_a^+\) be its projected endpoints. The reported genus-independent identity is
\[
 E(u)=\int u^*\omega
 =\epsilon\sum_a\bigl(f_{i(a)}(q_a^+)-f_{i(a)}(q_a^-)\bigr).
\]
For a fixed finite arc pattern and bounded endpoint actions, this gives \(E(u)=O(\epsilon)\). For example bounded functions on a compact base give a constant controlled by their oscillations and the number of arcs. It is not asserted uniform when those data grow with genus or when additional Hamiltonian/interior-end action terms are introduced.

Exactness forces compact closed bubble components to be constant. It also forces a bordered component to be constant when no boundary label changes occur, so each closed boundary component lies on one exact graph and contributes zero action. Consequently a nonconstant component in this setting requires label-changing action. These statements concern the unperturbed exact model, not every punctured-curve or perturbed Floer problem.

Independently, a genuine graph spine of a connected compact oriented surface of genus \(g\) with \(b\geq1\) boundary components has
\[
 b_1(\Gamma)=2g+b-1.
\]
The formula applies to a spine of that surface, not automatically to its visible nonconstant image. Interior punctures require corresponding topological adjustments. Handles can collapse into constant ghost components while the visible flow skeleton has fewer cycles, even becoming a tree. The aid therefore proposes a stable decorated Morse graph retaining genus, boundary and marking information at ghost vertices rather than an undecorated flow graph alone. Precise decorations and stability data are not supplied. These checkpoints do not prove higher-genus compactness, a moduli-space identification or a gluing theorem; they specify constraints a meaningful degeneration theory must retain.
\subsection{Short English statement}
The problem asks to extend the disc degeneration to higher genus. The supplied aid gives an exact action/energy identity independent of genus and explains why lost handles cannot simply disappear from the limiting data. Constant ghost components must retain topology in a stable decorated Morse-graph target; a full higher-genus correspondence is not supplied.
\subsection{Sources}
\begin{itemize}
\item Local source record: catalogue entry, catalogue entry; canonical \texttt{aim-geometry-notes.json}, \texttt{notes[332]}, item A.13, ``Moduli space of holomorphic polygons,'' contributed by Yong-Geun Oh. Title: ``Exactness and topology in the higher-genus adiabatic limit.''
\item AIM workshop ``Amoebas and tropical geometry,'' supplied reference \url{https://aimath.org/WWN/amoebas/amoebas.pdf}. This draft uses the supplied record without accessing a PDF.
\item Supplied disc-theorem citation: Fukaya and Oh, ``Zero loop open strings in the cotangent bundle and Morse homotopy,'' \emph{Asian Journal of Mathematics} 1 (1997), 96--180, as recorded. The exactness-topology checkpoint and decorated-graph proposal belong to research attempt 1.
\item Registry dated 14 August 2026: ``partially\_solved,'' ``PARTIAL-PROGRESS,'' statement status ``exact''; attempt partial, rigor ``proved,'' novelty confidence low. The machine-generated, non-peer-reviewed aid requires independent verification before citation.
\end{itemize}
\end{document}
